\section{Introduction}

We start with an informal introduction to quantum computing and lay some motivations.

Fault-tolerant quantum computing has the potential, through algorithms like Shor's factoring algorithm \cite{Shor}, to deliver possibly exponential speedups over classical computation. However, current Noisy Intermediate-Scale Quantum (NISQ) computers have too few qubits and too much noise to deliver practical benefit \cite{PreskillNISQ}.

Computation in the NISQ era is often modelled by a hybrid setup: on the quantum computer, a noisy initial state is prepared, some noisy gates are applied, and qubits are measured. A classical computer then processes these measurements, and can modify and repeat the quantum computation if necessary.

For the quantum computer, there exist a variety of computational models: The quantum Turing machine \cite{QuantumTM}, the quantum circuit model \cite{QuantumCircuit}, adiabatic quantum computation \cite{Adiabatic}, and measurement-based models \cite{MBQCOriginal}, among many others. Though these models are equivalent from a complexity viewpoint (informally, any problem that is efficiently solvable in one is efficiently solvable in the others) \cite{MBQCOriginal,AdiabaticEqCircuit,CircuitEqTM}, the underlying quantum properties used by each may differ significantly.

The measurement-based quantum computation (MBQC) model is of particular interest. This is because MBQC collapses any quantum circuit into measurements of a highly-entangled resource state that may be prepared in constant depth. This is quite shocking, as it means that any quantum computation may be done in constant depth. The cost of this model, however, is that measurements need to be adaptive. In other words, processing must be done on-the-fly to determine what qubits should be measured next and how.

Since measurement is irreversible, MBQC leads to a model of one-way quantum computation. In other words, as measurement cannot increase the entanglement of a system, all entanglement required must be initially present in the resource state. Informally then, entanglement is the resource consumed by this model of quantum computation. Thus, there is interest in determining what resource states are universal for quantum computation.

Graph states are a kind of resource state whose data corresponds to a simple, undirected graph. The comparability grid is a graph with vertices $V = \{(x,y) \mid x,y \in [n]\}$ and an edge between $(x_1,y_1)$ and $(x_2,y_2)$ if $x_1 \geq x_2$ and $y_1 \geq y_2$, excluding self-loops. It is an open problem whether the family of all comparability grids is universal for MBQC. There are implications for both possibilities. For example, if comparability grids are not universal, then unbounded rank width is not a sufficient condition for the universality of a family of graph states, since the family of all comparability grids have unbounded rank width.

In this thesis, we report on progress to prove and disprove the universality of comparability grids. Concretely, we give a set of conjectures for how one may be able to extract universal quantum computation using comparability grids as our resource state. In tandem to this, we outline various algorithms to verify small cases of these conjectures. 

Readers with no prior exposure to quantum computation are encouraged to start at Appendix \ref{appendix:background}. In Section \ref{sec:background}, we introduce the necessary notation and background so that in Section \ref{sec:mbqc} we may formally define the concepts in our problem statement. Then, Section \ref{sec:cgrid} recalls some results regarding comparability grids. These theoretical results inspire the algorithmic approaches in Section \ref{sec:algorihmic-approaches}, which in turn guide the results in Section \ref{sec:conclusions}.